\documentclass[aspectratio=169,11pt]{beamer}

% --- Basic packages ---
\usepackage{amsmath,amssymb,amsfonts}
\usepackage{bm}
\usepackage{graphicx}
\usepackage{booktabs}
\usepackage{threeparttable} % For the notes section
\usepackage{multirow}   % For multi-row cells
\usepackage{mathtools}
\usepackage{hyperref}
\usepackage[dvipsnames]{xcolor}
\usepackage{enumitem}
\usepackage{forest}
\definecolor{PastelPink}{HTML}{FFB7C5}    % Baby Pink
\definecolor{PastelBlue}{HTML}{AEC6CF}    % Baby Blue
\definecolor{MintGreen}{HTML}{B7E4C7}     % Soft Green
\definecolor{SoftYellow}{HTML}{FFFACD}    % Lemon Chiffon
\definecolor{SoftLavender}{HTML}{E6E6FA}  % Light Purple
\definecolor{PastelPeach}{HTML}{FFDAB9}   % Light Orange
\definecolor{Periwinkle}{HTML}{CCCCFF}    % Light Blue-Violet
\definecolor{LightCoral}{HTML}{F08080}    % Soft Red
% Define the bullet shapes globally for the entire document
\setbeamerfont{itemize/enumerate body}{size=\normalsize}

% level 2: nested \begin{itemize}
\setbeamerfont{itemize/enumerate subbody}{size=\small}

% level 3: nested inside that
\setbeamerfont{itemize/enumerate subsubbody}{size=\footnotesize}
 % First level (main item)
\setlist[itemize,1]{
    label=\textcolor{softaccent}{\large\textbullet},
    itemsep=7pt,          % space between items
    topsep=5pt            % space above/below the list
}

% Level 2 bullets: smaller spacing
\setlist[itemize,2]{
    label=\textcolor{softaccent}{\small\textendash},
    font=\footnotesize,
    itemsep=4pt,          % tighter spacing
    topsep=4pt
}
\setlist[itemize,3]{
    label=\textcolor{softaccent}{\small\circ},
    itemsep=2pt,          % tighter spacing
    topsep=2pt
}
% --- Beamer setup ---
\usetheme{default}
\usecolortheme{dolphin}

% Soft colors
\definecolor{softbg}{HTML}{F7F5F0}      % warm, light background
\definecolor{softtext}{HTML}{222222}    % dark gray instead of black
\definecolor{softaccent}{HTML}{3B6A8F}  % muted blue for structure

% Overall text & background
\setbeamercolor{background canvas}{bg=softbg}
\setbeamercolor{normal text}{fg=softtext,bg=softbg}

% Titles and structure
\setbeamercolor{title}{fg=softaccent}
\setbeamercolor{frametitle}{fg=softaccent}
\setbeamercolor{structure}{fg=softaccent}

\usepackage{sourcesanspro}
\renewcommand{\familydefault}{\sfdefault}

% Frame title font
\setbeamerfont{frametitle}{size=\Large}

% Templates
\setbeamertemplate{navigation symbols}{}
\setbeamertemplate{footline}[frame number]
\setbeamertemplate{itemize items}[circle]
\setbeamertemplate{section in toc}[sections numbered]

% --- Document Metadata ---
\title{Reshaping Adolescents’ Gender Attitudes: Evidence from a School-Based Experiment in India}

\author{Diva Dhar, Tarun Jain, and Seema Jayachandran}
\institute[Your University]{\large AER 2022}
\date{Presentation by Arjab Bose}



% --- Custom Commands (for math in the slides) ---
\newcommand{\qhi}{q_h}
\newcommand{\qlo}{q_l}
\newcommand{\rh}{\rho_h}
\newcommand{\rl}{\rho_l}
\newcommand{\muup}{\mu^s}
\newcommand{\muorig}{\mu}
\newcommand{\bup}{b^s}
\newcommand{\bg}{b^g}
\newcommand{\bb}{b^b}

\begin{document}

% =================================
% SLIDE 1: TITLE PAGE
% =================================
\begin{frame}
  \titlepage
\end{frame}

% =================================
% MOTIVATION & OVERVIEW
% =================================
\begin{frame}{Motivation}
  \begin{itemize}
    \item  \textcolor{RoyalBlue}{Gender inequality} is acute in many developing countries, and it is deeply rooted in cultural norms.
   
    
    
   
   \item India, and Haryana in particular, shows this clearly:
      \begin{itemize}
        \item Only 0.80 girls enroll in tertiary education for every boy. 
        \item Child marriage is common. 
        \item India has one of the world's lowest female labor force participation rates.
        \item The child sex ratio is 1.09 boys per girl nationally and 1.20 in Haryana, the most male-skewed state. 
      \end{itemize}
    
    
    \item \textcolor{RoyalBlue}{Key Question: Can a school-based discussion program make adolescents' gender attitudes and behavior more equal, and do the effects persist?}
    
  \end{itemize}
\end{frame}

\begin{frame}{Introduction}
  \begin{itemize}
    \item Past literature has focused on policy constraints. This paper tests whether discussion and persuasion can change \textcolor{RoyalBlue}{preferences}. 

    \item The intervention was a school-based program named Legion of Stars. 
   \item Method $\rightarrow$ \textcolor{RoyalBlue}{Randomized Control Trial (RCT)} 
      
    \item What did the program do?
   \begin{itemize}
       \item 45-minute classroom discussions every three weeks for 2.5 school years, for students in grades 7-10.
       \item Facts and endorsement of gender equality, reflection on their own and society's views (stereotypes, roles at home, girls' education, women's employment, harassment). 
       \item Some sessions taught communication skills, such as persuading parents to allow later marriage.
       
   \end{itemize}
      
     
     
\end{itemize} 
\end{frame}

\begin{frame}{Research Question \& Results Preview}
  \begin{itemize}
    \item \textcolor{RoyalBlue}{Can a school-based intervention that promotes gender equality change gender attitudes of adolescents?}
    \item \textbf{Main outcome:} Gender attitudes measured as an index of survey responses on women's roles and rights (e.g., women working outside the home, more women in politics).
   
    \item \textbf{Findings:}
    
      \begin{itemize}
        \item The program shifted attitudes toward gender equality by 0.18 standard deviations after 3.5 months. 
        \item Approximately 16\% of the initial gender-regressive views were converted into support for gender equality.
        \item Two years after the program ended, the effect was still 0.16 standard deviations.
       
        
      \end{itemize}
    
  \end{itemize}
\end{frame}

% =================================
% THEORY
% =================================
\begin{frame}{Data and Empirical Strategy}
  \begin{itemize}
    \item \textcolor{RoyalBlue}{Sample}: 314 government secondary schools in 4 Haryana districts (150 treatment, 164 control), about 14,000 students surveyed. Randomization was at the school level.
    \item \textcolor{RoyalBlue}{Surveys}: Baseline in 2013-14 (grades 6-7, about 12 years old), end line 1 about 3.5 months after the program ended, end line 2 about 2 years later.

    \item \textcolor{RoyalBlue}{Empirical Specification}: 

\begin{equation}
Y_{ij} = \beta_0 + \beta_1 \, Treated_j + \beta_2 \, Y^0_{ij} + \beta_3 \, X_{ij} + \epsilon_{ij}
\end{equation}

\begin{itemize}
    \item $Y_{ij}$: end line outcome for student $i$ in school $j$
    \item $Treated_j$: equals 1 if school $j$ was assigned to treatment, 0 otherwise
    \item $\beta_1$: Average treatment effect 
    \item $Y^0_{ij}$: baseline analog of the outcome
    \item $X_{ij}$: grade-gender and district-gender fixed effects
   
\end{itemize}


  \end{itemize}
\end{frame}




\begin{frame}{Effect on Gender Attitudes: Short Run vs. Medium Run}

\centering
\small
\newsavebox{\tablebox}
\sbox{\tablebox}{%
\begin{tabular}{lcc}
\hline\hline
 & \begin{tabular}[c]{@{}c@{}}Gender attitudes index\\ (3.5 months after)\end{tabular}
 & \begin{tabular}[c]{@{}c@{}}Gender attitudes index\\ (about 2 years after)\end{tabular} \\
 & (1) & (2) \\
\hline
Treated & 0.180 & 0.160 \\
 & [0.020] & [0.019] \\[4pt]
Control group mean & 0.000 & 0.333 \\
Basic controls & Yes & Yes \\
Number of students & 13,987 & 13,679 \\
\hline\hline
\end{tabular}%
}

\usebox{\tablebox}

\vspace{12pt}
\begin{minipage}{\wd\tablebox}
\footnotesize \raggedright
\textit{Notes:} Outcome is the gender attitudes index. All regressions control for the baseline analog of the outcome, grade-gender and district-gender fixed effects, and missing flags for each variable used to construct the index. Standard errors in brackets, clustered by school.
\end{minipage}

\end{frame}
% =================================
% DATA
% =================================
\begin{frame}{Coding/Workflow}
  \begin{itemize}
      \item This project uses Stata. 
      \item Everything runs from one script: Master do file $\rightarrow$ global directories $\rightarrow$ do files, ado files $\rightarrow$ work on data.
      \item Sequence: Baseline 1 cleaning $\rightarrow$ Endline 1 cleaning $\rightarrow$ Endline 2 cleaning $\rightarrow$ Analysis (Tables, Figures).  
      \item Key Packages: outreg2, egenmore, unique, estout/esttab, coefplot, corrtex, putexcel. \textcolor{RoyalBlue}{(ssc install)}
      \item \textcolor{RoyalBlue}{Key Features} : 
      \begin{itemize}
          \item Loops: automate repetitive tasks, making the code shorter and reproducible \textcolor{RoyalBlue}{(foreach)}. 
          
          \item Folder Structure: Data(Rawdata,Cleandata), Code(do files), Output(Figures,Tables,other results).
      \end{itemize} 
  \end{itemize}
\end{frame}

\begin{frame}{Folder Structure}

\centering

\begin{forest}
for tree={
    grow'=0,
    align=left,
    edge={->},
    parent anchor=east,
    child anchor=west,
    l sep=15pt,
    s sep=7pt,
    font=\large
}
[Research Folder
    [Data
        [Raw Data]
        [Clean Data]
    ]
    [Code
        [Do Files]
    ]
    [Output
        [Figures]
        [Tables]
        [Other Results]
    ]
]
\end{forest}

\end{frame}

% =================================
% EMPIRICS: PREDICTION 1
% =================================
\begin{frame}[fragile]{Coding/Special Requirements}
  \begin{itemize}
    \item \textcolor{RoyalBlue}{global macros}: the main purpose is to store a folder location once and then reuse it throughout.
    \item Can be called with \textcolor{RoyalBlue}{\$name}, e.g.:
  \end{itemize}

  \vspace{0.5pt}
  {\footnotesize
  \begin{verbatim}
global main_loc "F:\Exeter_coursework\...\149882-V1"
global deid "$main_loc/data"
global baseline_student_raw "$deid/baseline_student_raw"
  \end{verbatim}
  }
  \begin{itemize}
      \item Write comments to organize do file (*****).
      \item Requirements: 
      \begin{itemize}
          \item License to Stata or other software.
          \item Central High Performance Computing Cluster/Supercomputer.
          \item Understanding Stata is strictly \textcolor{RoyalBlue}{case sensitive}.
          
      \end{itemize}
  \end{itemize}
\end{frame}



% =================================
% EMPIRICS: PREDICTION 2
% =================================
\begin{frame}{README}
  \begin{itemize}
    
    \item \textcolor{RoyalBlue}{What the README does well:}
    \begin{itemize}
        \item Clearly identifies the \textbf{master do-file} and explains that it runs the analysis in sequence.
        \item Provides a detailed list of the \textbf{data, code files, software requirements ,and folder structure}.
        \item Clear Diagram of folder workflow. 
        
    \end{itemize}
    \item \textcolor{RoyalBlue}{What could be clearer:}
    \begin{itemize}
        
        \item A map in the paper requires \textbf{ArcGIS}, but the README does not specify anything on it or provide the underlying step-by-step workflow. 
    
  \end{itemize}
  \end{itemize}
\end{frame}


% =================================
% Conclusion
% =================================
\begin{frame}{Replication and Conclusion}
  \begin{itemize}
    \item The do file was not hard to replicate apart from minor issues such as fixing paths and installing packages. 
    \item The map (districts covering the survey) could not be reproduced. 
    \item The slides are prepared in Latex (Overleaf). 
    \item The results of the paper is important for public policy makers/governments to rethink their approach in schools in India. 
    \item \textcolor{RoyalBlue}{Takeaway}: Very organized and smooth do file setup. 


    \end{itemize}

    
    
  
\end{frame}





\end{document}

